\section{Arquitectura orientada a eventos en LinkedIn}

\paragraph{Descripción y escala.}
LinkedIn utiliza una arquitectura orientada a eventos para actualizar los datos
que alimentan su sistema de recomendaciones. Las acciones de los usuarios
se publican como eventos en Apache Kafka, y los flujos de Apache Beam los
procesan para actualizar en segundos las variables que usan los modelos de
recomendación. En 2023, LinkedIn reportó que su infraestructura de procesamiento
en conjunto manejaba más de 4 billones de eventos diarios a través de más de
3\,000 flujos de datos \cite{xia2023linkedin}.

\paragraph{Motivación y arquitectura anterior.}
Antes del cambio, LinkedIn calculaba por lotes las variables que alimentaban
sus modelos de recomendación. Entre una acción del usuario y su efecto en
las recomendaciones podían pasar entre 24 y 48 horas, lo que dificultaba
reflejar sus intereses más recientes. Este retraso motivó el cambio al
procesamiento por eventos para actualizar esas variables a medida que llegaban nuevas acciones,
sin esperar a la siguiente ejecución por lotes \cite{xia2023linkedin}.

\paragraph{Beneficios y desafíos.}
Con el nuevo flujo, los datos derivados de la actividad de los usuarios
quedaban disponibles para los modelos de recomendación en pocos segundos.
Un desafío fue que los equipos de IA tenían que integrar varias herramientas
y encargarse del despliegue y monitoreo de los flujos. LinkedIn abordó este
problema con Managed Beam, una plataforma interna que automatizó parte de
esas tareas y permitió a los equipos centrarse en definir qué datos necesitaban
sus modelos \cite{xia2023linkedin}.
